% Source: https://artofproblemsolving.com/community/c3h2591396_php_prob
% Stable identity: aops-topic-2591396-q01
Let $s_k$ be the number of chocolates eaten in the first $k$ days. Since he eats at least one per day,
\[1\le s_1<s_2<\cdots<s_{58}=100.\]
Consider the $116$ numbers
\[s_1,s_2,\dots,s_{58}\qquad\text{and}\qquad s_1+15,s_2+15,\dots,s_{58}+15.\]
They all lie between $1$ and $115$, so by the pigeonhole principle two of them are equal. The $s_k$ are all different, and so are the $s_k+15$, so we must have $s_j=s_i+15$ for some $i$ and $j$; then $j>i$. On days $i+1,\dots,j$ he ate $s_j-s_i=15$ chocolates. $\blacksquare$
